\hwhat{Hawkes fits to agent talk events in 35 non-holdout goal periods (282 days), each fitted separately with a schedule baseline. Predicted: shared-objective weeks run near critical (median $\hat n\ge0.7$) and above free weeks, and \#51 drifts toward $n=1$ as its roster grows. All failed. Median $\hat n$ is 0.41 (0.36 in round 1b), bracketed at 0.2--0.4 by the baseline choice. Shared-objective weeks gave 0.32 vs 0.53 for free weeks ($p=0.75$); $\hat n$ \emph{falls} with swarm size ($\rho=-0.53$), and in \#51 as its roster grows ($\rho=-0.64$). A constant baseline would have ``confirmed'' criticality ($\hat n\approx0.99$), and pooled $\hat n$ cannot be separated from 10--30-min rate modulation. The robust positive: fast cross-agent triggering, $n_\times\approx0.07$ at 10--30 s lags, about 4$\times$ an agent-shift null ($p=2\times10^{-9}$). Humans and the nudger directly drive about 1.5\% of events. \textbf{Round 1b} (corrected exogenous inputs, long days cut at operator-off gaps, day-edge-trimmed nulls): every period verdict is unchanged, fast cross-triggering survives the corrected null, and withdrawing all operator messages inside \#51 (NE43) leaves $\hat n$ and activity unchanged. No holdout run yet.}

\hmodel{Model 09 (Hawkes self-exciting point process). Events: agent messages; one realization per village day. $\mu_B(t)$: schedule baseline (per-day level $\times$ 30-min shape, plus a kickoff bump); $g_{\text{exo}}$: human and nudger input; $\phi$: triggering kernel; $n$: branching ratio, the mean number of events each event directly triggers (critical at $n=1$; $\langle s\rangle$: mean cascade size). Multivariate fit: $\lambda_a$, agent $a$'s intensity, with own and cross kernels; $\bar m$: active agents; $w^{\times}_m$: cross weight at timescale $\tau_m$.}

\hmath{\vspace{-5pt}\begin{gather*}
\lambda(t)=\mu_B(t)+\textstyle\sum_{x<t}g_{\text{exo}}(t-x)+\sum_{t_j<t}\phi(t-t_j)\\
\phi(\tau)=n\beta e^{-\beta\tau},\qquad \langle s\rangle=1/(1-n)\\
\lambda_a(t)=c_a\mu(t)+\textstyle\sum_{t_j\in a}\phi_{\text{self}}(t-t_j)+\sum_{t_j\notin a}\phi_{\times}(t-t_j)\\
n_\times=(\bar m-1)\textstyle\sum_{\tau_m\le300\,\text{s}}w^{\times}_m
\end{gather*}\vspace{-7pt}}

\hobs{../figures/summary_obs.pdf}{(a) Branching ratio per goal period against swarm size, by coupling mode, with day-bootstrap 95\% CIs: none is near critical, shared-objective periods sit no higher than free ones, and $\hat n$ falls with size. (b) Fast cross-agent excitation against an agent-shift null that breaks cross-agent timing: nearly every period sits above the diagonal, the one component that beats its strongest null.}

\htelos{A negative that matters a little: this swarm is nowhere near a runaway cascade, and neither more agents nor a shared goal pushes it there. Anyone worried about message storms gets mild reassurance, but activity is set mostly by the scheduler and roster, not by social contagion. The one robust number (about 0.07 extra messages from others within half a minute of each message) may just be agents waking when a message lands: a scaffold property. A naive constant-baseline Hawkes fit would have reported near-criticality. Practically: one pitfall and one small coupling constant.}

\hbottom{The swarm is clearly subcritical ($n\approx0.2$--$0.4$) and its self-excitation does not depend on the goal's coupling mode; only a small, fast, reply-like cross-triggering survives its null.}

\hresults{{\footnotesize\setlength{\tabcolsep}{2pt}\begin{tabular}{@{}p{0.38\columnwidth}p{0.47\columnwidth}p{0.12\columnwidth}@{}}
\textbf{Prediction} & \textbf{Observed (round 1 $\to$ round 1b)} & \textbf{Verdict}\\\hline
P1 free weeks subcritical ($\tilde n_F<0.5$) & 0.53 $\to$ 0.52; 4/8 $\ge0.5$ & $\times$\\
P2 shared $>$ free & 0.32 vs 0.53 $\to$ 0.29 vs 0.52 ($p=0.80$) & $\times$\\
P3 shared near critical ($\ge0.7$) & 0.32 $\to$ 0.29 & $\times$\\
P4 mode survives $N$, hours, regime & $t=-0.19\to-0.67$; $\log N$ $t=-2.06$ & $\times$\\
P5 \#51 drifts toward $n=1$ & $\rho=-0.64\to-0.66$ (own split), $-0.70$ (shared units) & $\times$\\
P6 fast $n_\times$: shared $>$ free & 0.075 vs 0.071 $\to$ 0.075 vs 0.073 & $\times$\\
Fast $n_\times$ vs agent-shift null & 0.070 vs 0.017 $\to$ 0.074 vs 0.011 after trimming ($p=2\times10^{-8}$) & $\checkmark$\\
Native NE43: drive withdrawn & $\hat n$ 0.44/0.35/0.47 (within 2 SD); talk rate $-11\%$ & $\checkmark$\\
Native \#2 / NE14 & $n_\times$ 0.16 vs null max 0.11; trim moves $\Delta\le0.03$, $\hat n_{\rm TALK}$ rises & $\checkmark$/$\sim$\\
\end{tabular}}}

\hobsb{../figures/summary_obs2.pdf}{(a) Per-period $\hat n$ in round 1 vs round 1b: corrected exogenous drive and day splitting move almost nothing; \#6 and \#38 change because a long operator-off gap is now cut out of one day. (b) NE43 inside \#51 at a fixed roster: $\hat n$ does not fall when the daily bookends stop (B) or when the nudger stops (C).}

\hcaveats{Pooled $\hat n$ is still not identified against 10--30-min rate modulation. In regime III, \texttt{events\_core} lacks computer-use turns, so ALL means something different. The M3 at-risk set ignores rooms. NE43's sides are confounded with the \#focus room and side C has 5 days. Round-1b bootstraps use 50/25 replicates (round 1: 100/50). In \#51, $N$ is collinear with date and drive. H42 is fitting the read-out kernel this round does not model.}

\hnext{Leave the kernel to H42; next for H03: an autonomy fraction from the context ledger (H03-R1), NE10 (first nudges), and locking the holdout predictions (fast $n_\times$ above the trimmed shift null; $\hat n<0.8$).}
